
\documentclass{article}
\usepackage{colortbl}
\usepackage{makecell}
\usepackage{multirow}
\usepackage{supertabular}

\begin{document}

\newcounter{utterance}

\centering \large Interaction Transcript for game `eqbench', experiment `eq\_bench', episode 69 with reflexion{-}Qwen3.5{-}27B.
\vspace{24pt}

{ \footnotesize  \setcounter{utterance}{1}
\setlength{\tabcolsep}{0pt}
\begin{supertabular}{c@{$\;$}|p{.15\linewidth}@{}p{.15\linewidth}p{.15\linewidth}p{.15\linewidth}p{.15\linewidth}p{.15\linewidth}}
    \# & \multicolumn{2}{c}{Player} && \multicolumn{2}{c}{Game Master} \\
    \hline

    \theutterance \stepcounter{utterance}  
    & & & \multicolumn{4}{p{0.6\linewidth}}{
        \cellcolor[rgb]{0.9,0.9,0.9}{
            \makecell[{{p{\linewidth}}}]{
                \texttt{\tiny{[P1$\langle$GM]}}
                \texttt{You have attempted this same task before and did not succeed. These are the plans you wrote after those attempts.} \\
\texttt{Trial \#1: New plan: My previous attempt failed because I likely misjudged the intensity of the emotions based on the specific context of the insult. While "Anger" was high, the specific nature of Olivia's final remark—mocking a culinary failure ("sprinkled cumin on your latest disaster")—is a direct attack on Jacques's professional pride and competence, which he previously claimed was superior. This specific type of public humiliation suggests "Embarrassment" should be scored significantly higher than I did, potentially rivaling or exceeding Anger, as the insult targets his core identity as a talented chef. Additionally, I should ensure that the scores reflect the cumulative effect of the entire exchange, not just the final line. Next time, I will re{-}evaluate the dialogue to identify the specific vulnerability attacked (in this case, professional skill) and assign a higher score to the emotion corresponding to that vulnerability (Embarrassment), while ensuring the sum of scores reflects the high tension of the scene. I will also double{-}check that at least one score is > 0, which I did, but the distribution was likely the error. I will adjust the plan to prioritize the emotion most directly triggered by the final, specific insult.} \\
\\ 
\texttt{Trial \#2: New plan: My previous attempt failed because I overestimated the intensity of "Embarrassment" relative to "Anger" in this specific context. While Olivia's insult about the "cumin" and "disaster" is humiliating, Jacques's entire dialogue is characterized by aggressive defensiveness, arrogance ("over{-}inflated ego"), and direct hostility ("peddling your half{-}baked opinions"). The final insult is a provocation that would primarily trigger a defensive, aggressive reaction (Anger) rather than a passive feeling of shame (Embarrassment). A character with such a large ego is more likely to react with rage to a challenge to their talent than to feel genuine embarrassment. Next time, I will prioritize "Anger" as the dominant emotion (score 9{-}10) because the dialogue is a heated argument where Jacques is actively attacking back, and I will lower "Embarrassment" (score 4{-}5) to reflect that while the insult is stinging, his primary reaction is aggression, not shame. I will ensure "Indifference" and "Fear" remain at 0 as they are inconsistent with the high{-}conflict tone.} \\
\\ 
\texttt{Trial \#3: New plan: My previous attempts failed because I oscillated between over{-}prioritizing "Embarrassment" (Trial 1) and over{-}prioritizing "Anger" (Trial 2) without a consistent analytical framework for the specific character archetype. Jacques is defined by an "over{-}inflated ego" and claims of superior talent. When such a character is publicly mocked regarding their core competence (the "cumin" insult), the psychological reaction is a volatile mix: the ego is bruised (causing deep Embarrassment), but the defense mechanism is immediate, aggressive hostility (causing high Anger). In a high{-}stakes verbal spar, these two emotions often peak simultaneously rather than one suppressing the other. My strategy will be to assign the maximum possible score (10) to both Anger and Embarrassment, acknowledging that the insult triggers both the feeling of shame and the rage to defend against it. I will keep Indifference and Fear at 0, as the character is actively engaged and aggressive. This dual{-}maximum approach accounts for the complexity of an ego{-}driven character facing a direct hit to their professional pride.} \\
\\ 
\\ 
\\ 
\texttt{Your task is to predict the likely emotional responses of a character in this dialogue:} \\
\\ 
\texttt{Jacques: Fancy seeing you here, Chase.} \\
\texttt{Olivia: The feeling is mutual, Duval. } \\
\texttt{Jacques: I see you're still peddling your half{-}baked opinions to anyone who'll listen.} \\
\texttt{Olivia: And I see you're still mistaking your over{-}inflated ego for talent.} \\
\texttt{Jacques: Talent? I've got more talent in my little finger than you have in your entire body.} \\
\texttt{Olivia: Oh, is that the same little finger that sprinkled cumin on your latest disaster of a dish?} \\
\texttt{{[}End dialogue{]}} \\
\\ 
\texttt{At the end of this dialogue, Jacques would feel...} \\
\texttt{Embarrassment} \\
\texttt{Anger} \\
\texttt{Indifference} \\
\texttt{Fear} \\
\\ 
\texttt{Give each of these possible emotions a score from 0{-}10 for the relative intensity that they are likely to be feeling each.} \\
\\ 
\texttt{You must output in the following format, including headings (of course, you should give your own scores), with no additional commentary:} \\
\\ 
\texttt{Embarrassment: <score>} \\
\texttt{Anger: <score>} \\
\texttt{Indifference: <score>} \\
\texttt{Fear: <score>} \\
\\ 
\\ 
\texttt{{[}End of answer{]}} \\
\\ 
\texttt{Remember: zero is a valid score, meaning they are likely not feeling that emotion. You must score at least one emotion > 0.} \\
\\ 
\texttt{Your answer:} \\
            }
        }
    }
     \\ \\

    \theutterance \stepcounter{utterance}  
    & \multicolumn{4}{p{0.6\linewidth}}{
        \cellcolor[rgb]{0.561,0.737,0.561}{
            \makecell[{{p{\linewidth}}}]{
                \texttt{\tiny{[P1$\rangle$GM]}}
                \texttt{Embarrassment: 9} \\
\texttt{Anger: 10} \\
\texttt{Indifference: 0} \\
\texttt{Fear: 0} \\
            }
        }
    }
    & & \\ \\

    \theutterance \stepcounter{utterance}  
    & & & \multicolumn{2}{p{0.3\linewidth}}{
        \cellcolor[rgb]{0.9,0.9,0.9}{
            \makecell[{{p{\linewidth}}}]{
                \texttt{\tiny{[GM$|$GM]}}
                \texttt{Embarrassment: 9} \\
\texttt{Anger: 10} \\
\texttt{Indifference: 0} \\
\texttt{Fear: 0} \\
            }
        }
    }
    & & \\ \\

    \theutterance \stepcounter{utterance}  
    & & & \multicolumn{2}{p{0.3\linewidth}}{
        \cellcolor[rgb]{0.9,0.9,0.9}{
            \makecell[{{p{\linewidth}}}]{
                \texttt{\tiny{[GM$|$GM]}}
                \texttt{Embarrassment: 7} \\
\texttt{Anger: 8} \\
\texttt{Indifference: 0} \\
\texttt{Fear: 0} \\
            }
        }
    }
    & & \\ \\

    \theutterance \stepcounter{utterance}  
    & & & \multicolumn{2}{p{0.3\linewidth}}{
        \cellcolor[rgb]{0.9,0.9,0.9}{
            \makecell[{{p{\linewidth}}}]{
                \texttt{\tiny{[GM$|$GM]}}
                \texttt{game\_result = LOSE} \\
            }
        }
    }
    & & \\ \\

    \theutterance \stepcounter{utterance}  
    & & \multicolumn{4}{p{0.6\linewidth}}{
        \cellcolor[rgb]{0.95,0.95,0.95}{
            \makecell[{{p{\linewidth}}}]{
                \texttt{\tiny{[P1$\leftrightarrow$P]}}
                \texttt{Post{-}episode Reflexion (not sent to GM):} \\
\texttt{New plan: My previous attempts failed because I treated the emotional response as a single dominant state, oscillating between prioritizing Anger or Embarrassment, rather than recognizing that the specific insult ("cumin on your latest disaster") targets Jacques's professional identity, which he explicitly claims is superior. For a character with an "over{-}inflated ego," a direct hit to their claimed talent triggers a simultaneous, intense reaction: deep shame (Embarrassment) at the exposure of incompetence, and immediate, defensive rage (Anger) to mask that shame. The dialogue shows a heated, escalating conflict where Jacques is actively attacking, ruling out Indifference and Fear. Therefore, the correct strategy is to assign the maximum score (10) to both Anger and Embarrassment, acknowledging that the ego{-}driven character experiences the sting of humiliation and the fire of retaliation with equal intensity. I will output scores of 10 for both Anger and Embarrassment, and 0 for Indifference and Fear.} \\
            }
        }
    }
    & \\ \\

\end{supertabular}
}

\end{document}
